\documentclass[11pt,a4paper]{article}
\usepackage[T1]{fontenc}
\usepackage[utf8]{inputenc}
\usepackage{lmodern}
\usepackage[margin=2.2cm]{geometry}
\usepackage{amsmath,amssymb}
\usepackage{listings}
\usepackage{xcolor}
\usepackage{booktabs}
\usepackage{enumitem}
\usepackage{hyperref}
\hypersetup{colorlinks,linkcolor=black,urlcolor=blue!50!black}

% notation of the TEQ paper
\newcommand{\bl}[1]{\lfloor #1\rfloor}                       % butlast
\newcommand{\tr}[3]{\mathsf{trace}(#1,\ \mathsf{for}\ #2\ \mathsf{do}\ #3)}
\newcommand{\mapc}[2]{\mathsf{maps}(\mathsf{for}\ #1\ \mathsf{yield}\ #2)}
\newcommand{\al}[2]{\mathsf{invs}(\mathsf{for}\ #1\ \mathsf{check}\ #2)}
\newcommand{\cnt}{\mathsf{counter}}
\newcommand{\cur}{\mathsf{current}}
\newcommand{\ear}{\mathsf{earlier}}
\newcommand{\intf}{\mathsf{integer}}
\newcommand{\natf}{\mathsf{natural}}
\newcommand{\stores}{\mathsf{stores}}
\newcommand{\selects}{\mathsf{selects}}
\newcommand{\consts}{\mathsf{consts}}
\newcommand{\occurs}{\mathsf{occurs}}
\newcommand{\indexof}{\mathsf{indexof}}
\newcommand{\prefix}{\mathsf{prefix}}
\newcommand{\ite}[3]{\mathsf{if}\ #1\ \mathsf{then}\ #2\ \mathsf{else}\ #3}
\newcommand{\lemma}[1]{\textsf{#1}}
\newcommand{\sel}[2]{#2.#1}                                   % component #1 of state #2
\newcommand{\rw}{\ \leadsto\ }

\lstdefinelanguage{dafny}{
  morekeywords={method,returns,requires,ensures,var,while,if,else,function,lemma,forall,int,bool,array,datatype,case,for,do,yield,check,trace,maps,invs},
  sensitive=true, morecomment=[l]{//}, morestring=[b]",
}
\lstset{language=dafny,basicstyle=\ttfamily\footnotesize,keywordstyle=\bfseries,commentstyle=\itshape\color{black!60},
  columns=fullflexible,keepspaces=true,xleftmargin=1.5em}

\title{Verifying Loops without Invariants:\\ Trace Equation Systems by Two Examples}
\author{Gidon Ernst\\[2pt] \small LMU Munich, Germany\\ \small\texttt{gidon.ernst@gmail.com}}
\date{October 13, 2026}

\begin{document}
\maketitle

\begin{abstract}
This note explains how the verifier Cuvée proves array programs without loop invariants, by two examples that
every reader of a verification paper knows: the initialization of an array and the element-wise exchange of two
arrays. A loop is represented as a fold over its iterations whose result is a \emph{trace}, the sequence of values of
the program state at the loop head, and the verification condition is a first-order formula over traces. The formula
is proved by rewriting with a database of lemmas about traces, each a small fact about a loop shape (a counter, a
sequence of stores, a constant), under a context that keeps the facts of the goal. No invariant is written, and no
lemma is written for the two programs. The note introduces the representation, the lemma database, and the
engine as the examples need them, and ends with what the examples do not show.
\end{abstract}

\section{The idea}

A verifier based on Hoare logic asks for an invariant of every loop. For the initialization of an array the invariant
is easy, \lstinline|forall k :: 0 <= k < i ==> a[k] == v|, but it is an input of the user and not a property of the
program, and for programs that compute something more involved it is where the effort goes. Trace equation systems
take another route. Instead of asking what holds at the loop head in every iteration, they ask what the loop
\emph{computes}: the values of the variables after $n$ iterations, as a function of the inputs and $n$. For a loop
that counts, that function is $i_0 + n$; for a loop that stores $v$ at positions $0, \ldots, n-1$, it is the array
with those positions set to $v$. Once the loop is summarized by such a function, the postcondition is a claim about
the function, and that claim is first-order.

The summary is obtained in two steps. Symbolic execution turns the loop into a \emph{trace}, a sequence of states,
defined recursively from the body: the state after an iteration is the body applied to the state before it. This is
mechanical and needs no insight. Then a rewrite engine applies lemmas from a database to the trace, each stating what a
particular shape of loop computes, until the trace has a closed form and the postcondition follows by arithmetic.
The insight sits in the database, written once; a new program is verified without a lemma of its own as long as its
loops are composed of shapes the database knows.

The two examples below are such programs. The first shows the representation and the ordinary lemmas, the
second shows what happens when two variables of a loop depend on each other, which is where the approach needed
an idea that the paper on trace equation systems did not have.

\section{Array initialization}

\begin{lstlisting}
method init(n: int, v: int) returns (a: array<int, int>)
  requires 0 <= n
  ensures forall k: int :: 0 <= k < n ==> a[k] == v
{
  var i: int;
  i := 0;
  while i < n
  {
    a[i] := v;
    i := i + 1;
  }
}
\end{lstlisting}

Arrays here are total maps from integers to integers (the type \lstinline|array<int, int>|), so there is no bounds
check and \lstinline|a[i] := v| is the functional update $a[i := v]$. The run is
\begin{center}
\texttt{./Cuvee.sh -teq-state -prove examples/dafny/state/init.dfy}
\end{center}
where \texttt{-teq-state} produces the verification condition over traces and \texttt{-prove} runs the engine. The
output is empty: the goal is proved and therefore dropped.

\subsection{The verification condition}

Symbolic execution runs the method body on symbolic values, collecting the assumptions and checks of each path. The
loop is the only interesting statement. Its modified variables are $a$ and $i$, so the loop's state is the pair
$(a, i)$, and the condition produced for the method is (with \texttt{-print dfy} in place of \texttt{-prove}, the
fresh variables numbered by the tool)
\begin{align*}
&\forall n, v, a.\ 0 \le n \Rightarrow \forall a', i', n_4{:}\mathsf{Nat}, K, S.\\
&\quad K = \cnt(0, 1, n_4)\\
&\quad \wedge\ S = \tr{(a, 0)}{\bl{K}}{s, k :: (\sel{0}{s}[\sel{1}{s} := v],\ \sel{1}{s} + 1)}\\
&\quad \wedge\ a' = \sel{0}{\cur(S)}\ \wedge\ i' = \sel{1}{\cur(S)}\\
&\quad \wedge\ \al{\bl{K}, \bl{S}}{k, s :: \sel{1}{s} < n}\\
&\quad \wedge\ 0 \le \intf(n_4)\\
&\quad \Rightarrow \neg(i' < n) \Rightarrow \forall k.\ 0 \le k < n \Rightarrow a'[k] = v
\end{align*}
Read it line by line; every construct of the approach appears once.

\paragraph{The number of iterations.} $n_4$ is a natural number, the number of times the body runs. Nothing is
known about it except that it is some natural (its integer value $\intf(n_4)$ is nonnegative, stated explicitly
since the arithmetic works on integers). The loop test does not determine it directly; it is constrained by the
iteration constraint below.

\paragraph{The counter.} $K = \cnt(0, 1, n_4)$ is the sequence $0, 1, \ldots, n_4$, a trace of $n_4 + 1$ integers:
the iteration numbers. Every loop gets such a counter, and the lemmas of the database are indexed by it.

\paragraph{Traces.} A trace of type $\mathsf{Trace}\langle T\rangle$ is a nonempty sequence of values of type $T$,
the values of some quantity at the loop head before each iteration and after the last. Two functions take it
apart: $\bl{X}$ (``butlast'', the floor brackets) is the list of all but the last element, i.e.\ the values before
each of the $n$ iterations, and $\cur(X)$ is the last element, the value after the loop. The $j$-th element of a
list is $\ear(j, \bl{X})$, with $j$ a natural number.

\paragraph{The state trace.} $S$ is the trace of the loop's state, a pair, written with the combinator
\[
\tr{s_0}{\bl{K}}{s, k :: u}
\]
It starts with $s_0 = (a, 0)$, the state before the first iteration, and for each
iteration $k$ of the counter its next value is the update $u$ applied to the previous value $s$. Here the update is
the loop body as a function of the state: the array component gets $v$ stored at the index component, the index is
incremented. The components of a pair are written $\sel{0}{s}$ and $\sel{1}{s}$ (the tool prints them as the
selectors \texttt{tuple2\_0(s)}, \texttt{tuple2\_1(s)} of the pair datatype of its Dafny prelude). The combinator is
a fold: $S$ has one element more than $\bl{K}$, and
\[
\cur(\tr{s_0}{\bl{K}}{s, k :: u}) = u(\cdots u(u(s_0, 0), 1) \cdots, n_4 - 1).
\]

\paragraph{The exit values.} After the loop, the program variables are the components of the final state,
$a' = \sel{0}{\cur(S)}$ and $i' = \sel{1}{\cur(S)}$, and the loop was left because its test failed, $\neg(i' < n)$.

\paragraph{The iteration constraint.} $\al{\bl{K}, \bl{S}}{k, s :: \sel{1}{s} < n}$ states that the loop test
held before every iteration: the combinator $\mathsf{invs}$ ranges over the two lists $\bl{K}$ and $\bl{S}$ in
parallel, binding the iteration number $k$ and the state $s$ at position $k$, and requires its body for all
positions. This is the only place where the number of iterations is pinned down: the test must hold $n_4$ times and
fail afterwards.

Nothing in this condition is specific to the program but the update and the test; the rest is the schema that every
loop gets. In particular there is no invariant.

\subsection{The lemma database}

The engine does not know what a counter or a store is; it knows rewrite rules. A rule is a lemma with the attribute
\lstinline|{:rw}| whose conclusion is an equation $l = r$ (or an equivalence), read from left to right, possibly
with preconditions as side conditions. The variables of the lemma are matched against the goal. The database of
\texttt{-teq-state} is one Dafny file, \texttt{resources/cuvee/teq/state.dfy}, with the definitions of the types and
functions above and about forty lemmas. The ones this example needs:
\begin{lstlisting}
lemma {:rw} counter_def(i: int, c: int, n: Nat)
  ensures trace(i, for butlast(counter(0, 1, n)) do x, k :: x + c) == counter(i, c, n)

lemma {:rw} current_counter(i: int, c: int, n: Nat)
  ensures current(counter(i, c, n)) == i + integer(n) * c

lemma {:rw} earlier_counter(k: Nat, i: int, c: int, n: Nat)
  ensures earlier(k, butlast(counter(i, c, n))) == i + integer(k) * c

lemma {:rw} invs_counter_upper_nonnegative(m: int, n: Nat)
  requires 0 <= m
  ensures invs(for butlast(counter(0, 1, n)) check k :: k < m) <==> integer(n) <= m

function {:fold} stores<a, b>(a: array<a, b>, keys: List<a>, values: List<b>): Trace<array<a, b>>
{ trace(a, for keys, values do b, k, v :: b[k := v]) }

lemma {:rw} maps_const<a>(x: a, i: int, c: int, n: Nat)
  ensures maps(for butlast(counter(i, c, n)) yield v :: x) == butlast(consts(n, x))

lemma {:rw} current_stores<a, b>(k: a, a: array<a, b>, keys: List<a>, values: List<b>)
  ensures current(stores(a, keys, values))[k]
    == if occurs(k, keys) then earlier(indexof(k, keys), values) else a[k]

lemma {:rw} occurs_counter_up(j: int, i: int, n: Nat)
  ensures occurs(j, butlast(counter(i, 1, n))) <==> i <= j < i + integer(n)

lemma {:rw} indexof_counter_up(j: int, i: int, n: Nat)
  requires i <= j < i + integer(n)
  ensures indexof(j, butlast(counter(i, 1, n))) == natural(j - i)

lemma {:rw} earlier_consts<a>(k: Nat, n: Nat, x: a)
  ensures earlier(k, butlast(consts(n, x))) == x
\end{lstlisting}
Three kinds of lemma appear. \emph{Introduction} lemmas recognize a loop shape and name it: a trace that adds a
constant is a counter (\lemma{counter\_def}); a trace that stores values at keys is $\stores(a, keys, values)$,
which is a function definition marked \lstinline|{:fold}|, i.e.\ its body is rewritten \emph{to} the function.
\emph{Acceleration} lemmas give the closed form of a named shape: the final value of a counter, its $k$-th value,
the element of a store sequence at a key (\lemma{current\_stores}: the last value stored there, or the original
if the key was never stored). \emph{Bookkeeping} lemmas relate the list functions: $\occurs$, $\indexof$, $\prefix$,
the constant list $\consts$, and the third combinator $\mapc{xs}{x :: e}$, the list of $e$ for each element of $xs$,
which the engine introduces itself when a lemma needs a sequence of values that the program computes implicitly.

These lemmas are proved once, in Lean, from the recursive definitions; the program proof trusts them.

\subsection{The proof}

The engine simplifies the goal under a context of facts, innermost subterms first, and after each round applies the
first rule whose left-hand side matches and whose side conditions hold. Every rewrite step can be printed
(\texttt{-debug rewrite}); what follows is that printout, cleaned up. Three ideas appear along the way that are not
in the condition itself: contexts of combinators, projections, and views.

\paragraph{The index is the iteration number.} The engine first simplifies the update of $S$. To do so it builds the
\emph{context of the combinator}: inside $\mathsf{for}\ \bl{K}, \bl{S}\ \mathsf{do}\ s, k$, the bound variable $k$
is a value of $\bl{K}$, so $0 \le k < \intf(n_4)$, and $s$ is the state at position $k$, $s = \ear(\natf(k), \bl{S})$;
these are facts of the context in which the update is simplified. The update reads $\sel{1}{s}$, the index
component. The engine never writes down a trace for the index; instead, when it meets the read $\sel{1}{s}$ of a
component, it computes the \emph{view} of that read: the trace of this component alone, the \emph{projection} of
$S$ to component~1,
\[
\pi_1(S) = \tr{0}{\bl{K}, \bl{S}}{c, k, s :: c + 1}
\]
obtained from $S$ by taking the first component of the initial state and of each result, and by replacing the read
of the own component by the accumulator $c$ (the other components are still read from $s$, hence $\bl{S}$ is a
dependency). This update does not use $s$, so the unused dependency is dropped and $\lemma{counter\_def}$ fires:
\begin{align*}
\tr{0}{\bl{K}}{c, k :: c + 1} &\rw \cnt(0, 1, n_4) \\
\ear(\natf(k), \bl{\cnt(0, 1, n_4)}) &\rw 0 + \intf(\natf(k)) \cdot 1 \rw k
\end{align*}
The second line is the value of the view at position $k$ by \lemma{earlier\_counter}, and $\intf(\natf(k)) = k$
since $0 \le k$ is known. The view is a closed term that reads no component any more, so the engine replaces the
read in place: in every iteration $k$, the index is $k$. This is the step where ``the loop stores at $i$'' becomes
``the loop stores at $0, 1, \ldots$''.

\paragraph{The array is a sequence of stores.} The update of $S$ now has the array component $\sel{0}{s}[k := v]$,
and the projection of $S$ to component~0 is $\tr{a}{\bl{K}}{c, k :: c[k := v]}$ (the index read is gone, so
$\bl{S}$ is dropped again). It is matched by the body of $\stores$ with $keys = \bl{K}$ and the values a sequence that
the program never names: the engine synthesizes it as the map of the stored expression over the iterations,
\begin{align*}
\tr{a}{\bl{K}}{c, k :: c[k := v]} &\rw \stores(a, \bl{K}, \mapc{\bl{K}}{k :: v}) \\
\mapc{\bl{K}}{k :: v} &\rw \bl{\consts(n_4, v)}
\end{align*}
by the fold and \lemma{maps\_const}. The view of the exit value $a' = \sel{0}{\cur(S)}$ is therefore
\[
\cur(\stores(a, \bl{K}, \bl{\consts(n_4, v)})),
\]
the array after storing $v$ at each key of the counter. The view of the final index is
\[
\cur(\cnt(0, 1, n_4)) \rw \intf(n_4)
\]
by \lemma{current\_counter}.

\paragraph{The number of iterations is $n$.} The iteration constraint has become
$\al{\bl{K}}{k :: k < n}$ after the index read was evaluated, and \lemma{invs\_counter\_upper\_nonnegative} turns
it into $\intf(n_4) \le n$ (its side condition $0 \le n$ is the precondition). The exit $\neg(i' < n)$ with
$i' = \intf(n_4)$ gives $n \le \intf(n_4)$, and the arithmetic of the context concludes $\intf(n_4) = n$. The engine
records this as $n_4 = \natf(i')$ and replaces $\intf(n_4)$ by $i'$ from here on (the rewrite \lemma{integer\_eq}
in the log).

\paragraph{The postcondition.} Finally $a'[k]$ for $0 \le k < n$: $a'$ is a variable whose definition in the context
is the marker $\sel{0}{\cur(S)}$, whose view is the store sequence, and \lemma{current\_stores} fires through both,
\begin{align*}
a'[k] &\rw \mathsf{if}\ \occurs(k, \bl{K})\ \mathsf{then}\ \ear(\indexof(k, \bl{K}), \bl{\consts(n_4, v)}) \\
&\qquad\ \mathsf{else}\ a[k] \\
\occurs(k, \bl{K}) &\rw 0 \le k < 0 + \intf(n_4) \qquad\text{true, since } \intf(n_4) = n \\
\indexof(k, \bl{K}) &\rw \natf(k - 0) \\
\ear(\natf(k), \bl{\consts(n_4, v)}) &\rw v
\end{align*}
by \lemma{occurs\_counter\_up}, \lemma{indexof\_counter\_up} and \lemma{earlier\_consts}. The conclusion $a'[k] = v$
is true, the goal is dropped. The whole run takes about two seconds, of which the JVM start is most.

The proof needed no fact about the loop except what the database knows about counters, stores and constant lists.
Had the program stored $v$ at $2i$ or counted down, other lemmas of the same database would have fired.

\section{The element-wise exchange of two arrays}

\begin{lstlisting}
method swap(a: array<int, int>, b: array<int, int>, n: int)
  returns (a_: array<int, int>, b_: array<int, int>)
  requires 0 <= n
  ensures forall k: int :: 0 <= k < n ==> a_[k] == b[k] && b_[k] == a[k]
{
  var i: int;
  var t: int;
  a_ := a;
  b_ := b;
  i := 0;
  while i < n
  {
    t := a_[i];
    a_[i] := b_[i];
    b_[i] := t;
    i := i + 1;
  }
}
\end{lstlisting}

The state of the loop now has four components, $(t, a', b', i)$, and the update of $S$ is
\[
s, k :: \big(\sel{1}{s}[\sel{3}{s}],\ \sel{1}{s}[\sel{3}{s} := \sel{2}{s}[\sel{3}{s}]],\
\sel{2}{s}[\sel{3}{s} := \sel{1}{s}[\sel{3}{s}]],\ \sel{3}{s} + 1\big).
\]
The rest of the condition has the same shape as before. What is new is that the two arrays read each other: the new
$a'$ stores a value of the old $b'$, and the new $b'$ a value of the old $a'$ (through $t$).

\subsection{Why this is the hard case}

In the paper on trace equation systems, a loop yields one trace per variable, and a variable that reads another has
the history of that other among its dependencies: $A' = \tr{a}{\bl{K}, \bl{B'}}{x, k, y :: x[k := y[k]]}$ and
symmetrically $B'$ over $\bl{A'}$. The two equations are recursive through each other. To summarize $A'$, the
engine needs $B'$ in closed form, which needs $A'$ in closed form. The paper resolves such cases by lemmas that tie
the two equations together, written per shape, and the summaries of the sorting algorithms in that line of work
carry a dozen side conditions each for this reason.

The state trace sidesteps the recursion: there is one fold over the pair of arrays, and its components are
projections computed on demand. The question is only how a projection may read the other component.

\subsection{The proof}

\paragraph{The index, as before.} The read $\sel{3}{s}$ becomes $k$ exactly as in the first example; afterwards
every store and read in the update is at position $k$.

\paragraph{The view of $a'$ reads $b'$.} To summarize the array $a'$, the engine computes its projection,
\[
\pi_1(S) = \tr{a}{\bl{K}, \bl{S}}{c, k, s :: c[k := \sel{2}{s}[k]]},
\]
which still reads component~2 of the state, so $\bl{S}$ stays a dependency. The fold of $\stores$ matches it with
the synthesized value sequence,
\[
\pi_1(S) \rw \stores(a, \bl{K}, \mapc{\bl{K}, \bl{S}}{k, s :: \sel{2}{s}[k]}),
\]
but that is not closed: the values stored are reads of $b'$. The engine now needs to know what $\sel{2}{s}[k]$ is,
the element $k$ of the other array in iteration $k$. It computes the view of the read $\sel{2}{s}$: the projection
of $S$ to component~2 at position $k$,
\[
\ear(\natf(k), \bl{\stores(b, \bl{K}, \mapc{\bl{K}, \bl{S}}{k', s' :: \sel{1}{s'}[k']})}),
\]
whose update in turn reads component~1. Here the engine would loop forever if it followed every read: the view of
$a'$ needs the view of $b'$, which needs the view of $a'$. The rule is that \emph{a component is not projected
inside its own view}: while the view of component~1 is being computed, a read of component~1 stays as it is, a
marker that the lemmas treat as an opaque value. So $\sel{1}{s'}[k']$ stays, and the view of $b'$ is a sequence of
stores whose values are unknown terms. That is enough, because of what is asked next.

\paragraph{Iteration $k$ reads a position that no earlier iteration wrote.} The lemma on an element of an earlier
value of a store sequence,
\begin{lstlisting}
lemma {:rw} earlier_stores<a, b>(k: a, j: Nat, a: array<a, b>, keys: List<a>, values: List<b>)
  ensures earlier(j, butlast(stores(a, keys, values)))[k]
    == if occurs(k, prefix(j, keys)) then earlier(indexof(k, prefix(j, keys)), values) else a[k]
\end{lstlisting}
says: the value at $k$ after $j$ stores is the last value stored at $k$ among the first $j$ keys, or the original if
$k$ is not among them. Applied to the read, with $j = \natf(k)$ and $keys = \bl{K}$,
\begin{align*}
\sel{2}{s}[k] &\rw \ite{\occurs(k, \prefix(\natf(k), \bl{K}))}{\ear(\ldots)}{b[k]} \\
\prefix(\natf(k), \bl{K}) &\rw \bl{\cnt(0, 1, \natf(k))} \\
\occurs(k, \bl{\cnt(0, 1, \natf(k))}) &\rw 0 \le k < 0 + \intf(\natf(k)) \qquad\text{i.e. } k < k,\ \text{false}
\end{align*}
by \lemma{prefix\_counter} and \lemma{occurs\_counter\_up}. The first $k$ keys are $0, \ldots, k-1$, and $k$ is not
among them, so the \textsf{then} branch, and with it the opaque read of $a'$, disappears: in iteration $k$ the loop
reads the \emph{original} $b[k]$. This is the whole argument of the program's correctness, and the engine found it by
asking the right question (what is the element at $k$ of the earlier value) rather than by unfolding both arrays.

\paragraph{The closed forms.} With the read resolved, the projection of $a'$ no longer depends on the state, and
\begin{align*}
\pi_1(S) &\rw \stores(a, \bl{K}, \mapc{\bl{K}}{k :: b[k]}) \rw \stores(a, \bl{K}, \selects(b, \bl{K})) \\
\pi_2(S) &\rw \stores(b, \bl{K}, \selects(a, \bl{K}))
\end{align*}
where $\selects(b, keys)$, the reads of $b$ at the keys, is another \lstinline|{:fold}| of the database. The exit
values are $\cur$ of these; the number of iterations is $n$ as before.

\paragraph{The postcondition.} $a'[k]$ for $0 \le k < n$ goes through the definition of $a'$, the view, and
\lemma{current\_stores} as in the first example, now with the values $\selects(b, \bl{K})$:
\begin{align*}
a'[k] &\rw \mathsf{if}\ \occurs(k, \bl{K})\ \mathsf{then}\ \ear(\indexof(k, \bl{K}), \selects(b, \bl{K}))\ \mathsf{else}\ a[k] \\
&\rw \ear(\natf(k), \selects(b, \bl{K})) \rw b[\ear(\natf(k), \bl{K})] \rw b[k]
\end{align*}
by \lemma{earlier\_selects} and \lemma{earlier\_counter}; $b'[k] = a[k]$ is symmetric. About three seconds in all,
fourteen distinct lemmas, none written for the program.

\section{The engine, in short}

The examples used the following machinery, which the engine of \texttt{-prove} provides for every goal.

\begin{description}[leftmargin=1.5em]
\item[Contextual simplification.] The goal is traversed innermost first; at each position the context holds the facts
  that are known there (the antecedents of the implication, the other conjuncts, the guards of a branch), and a
  subformula known to hold or to fail is replaced by its truth value. Combinator bodies get the context of their
  position: the ranges of the bound variables and their values in the dependencies.
\item[Rewriting with the database.] After each round, the first rule whose left-hand side matches the term at the
  top and whose side conditions hold is applied. Matching is modulo the usual algebra of arithmetic, the binding
  structure of the combinators, and higher-order patterns for shapes with a free function in them. Side
  conditions that are closed after the match are looked up among the facts and otherwise decided by the same
  simplification; conditions with free variables are matched against the facts, which binds them.
\item[Definitions and classes.] Equations among the facts, $K = \cnt(0, 1, n_4)$ or $a' = \sel{0}{\cur(S)}$, are kept
  in congruence classes. A variable that a pattern does not match is matched like each member of its class, so
  a lemma on counters applies to $\bl{K}$, and a lemma on stores applies to $a'$ through its definition and the view
  of that definition.
\item[Arithmetic.] Linear facts over the integers go to a simplex over the rationals with integer rounding, which
  decides the comparisons the examples needed ($\intf(n_4) = n$ from the two bounds, $k < k$ false).
\item[Views and suspension.] A read of a component of a state trace is matched modulo its view, the simplified
  projection of the state trace to that component; a component is not projected inside its own view; a read whose
  view is closed is replaced by it in place. Views are memoized by the facts that mention their variables.
\end{description}

The same two examples are proved by the earlier form of the tool (\texttt{-teq -simplify}, one trace per variable)
with the same lemmas; the swap needed there what the state form avoids, the two equations over each other's
histories, resolved by following the equations through the context. The state form matters for loops where
the variables cannot be layered, and most of all for nested loops, where an inner loop becomes a single value of the
outer update rather than a family of fresh sequences; those examples are in the same directory of the repository
(\texttt{examples/dafny/state}, nested sums, row sums, a transpose), and a binary search where the remaining gap is
visible.

\section{What the examples do not show}

Both programs have a layered loop: the index depends on nothing, the arrays depend on the index and on each other
only at the position being written. The database covers such loops without help. Programs whose loops carry an
accumulated property (a sorted prefix, a running maximum, a search that stops) need lemmas about their shapes, some
of them written for the program, and the engine has to match those lemmas against views whose shape it has not seen
before. The open items at the time of writing, with the binary search as the example: facts about a named trace
should become available when its term occurs, not only when the goal mentions the exact atom (the bounds of a
monotone search index), and lemmas over another dependency must be polymorphic in that dependency's type, since in
the state form it is the list of states. The record of the implementation,
\texttt{ideas/deduction-engine-implementation.md} in the repository, keeps the state of these.

\section*{AI Declaration}

This note was written with the assistance of Claude (Anthropic) in Claude Code, from the runs of the tool on the
two examples (\texttt{-debug rewrite}) and the documentation of the repository; the author directed the work and
is responsible for the content.

\end{document}
